% RESEARCH_DRAFT_SHA256: c3301c25e34b990cc3b3d03d3a7c0bdadf456e91a05072f90745a5015f5375eb
% Layout and source-reference wrappers expanded for the web renderer.
\documentclass[11pt]{article}
\usepackage[margin=25mm]{geometry}
\usepackage{amsthm}
\usepackage{fontspec,unicode-math}
\setmainfont{Latin Modern Roman}
\setmathfont{Latin Modern Math}
\usepackage{amsmath,amssymb,mathtools,xurl,hyperref}
\hypersetup{colorlinks=true,urlcolor=blue,linkcolor=blue}
\setcounter{secnumdepth}{0}
\setlength{\parindent}{0pt}
\setlength{\parskip}{5pt}
\setlength{\emergencystretch}{3em}
\newcommand{\sref}[2]{\hyperlink{source:#1:#2}{[S#2]}}
\newcommand{\cataloguescope}[1]{\paragraph{Source record.}#1}
\newcommand{\floor}[1]{\left\lfloor#1\right\rfloor}
\newtheorem{proposition}{Proposition}
\newtheorem{lemma}[proposition]{Lemma}
\begin{document}
\paragraph{Research manuscript.} OpenAI \texttt{gpt-6-astra}, reasoning \texttt{ultra}; 11 October 2026.
\paragraph{Preserved source body.} The complete catalogue mathematical text and sources follow with layout and source links expanded, including source qualifications and its historical preparation notes. New research begins at the section headed Research attempt.

% UnsolvedMath ID 30006672; source catalogue code OWR-14299916-010
\section{New Theoretical Principles for Foundation Models}\label{30006672}
{\small UnsolvedMath ID 30006672 · Computer Science}
\subsection{Definitions and mathematical statement}

\paragraph{Shared notation.}Standard number systems are $\mathbb N,\mathbb Z,\mathbb Q,\mathbb R,\mathbb C$. Specialized notation and index ranges follow the cited source and the definitions below.


Approximation theory studies representable functions and approximation error; optimization studies finding good parameters; statistical learning studies data-dependent prediction and its population loss.
\paragraph{Record scope and qualifications.}
The catalogue title is retained as metadata. Named certificates, contracts, bounds or frameworks in that title are not asserted as results of the original source extract.
Whether fundamentally new principles are needed is a conceptual research direction rather than a definite proposition.

Source-record limitation: Whether new theoretical principles are needed is a conceptual research direction, not a mathematical proposition.
\paragraph{Statement recorded in the supplied source.}
Do the phenomena observed in contemporary foundation models require fundamentally new mathematical principles, or can they be explained by extensions of existing theories of approximation, optimization, and statistical learning? \href{https://doi.org/10.4171/owr/2026/11}{[S2]}
\paragraph{Source status.}
The supplied catalogue labels this source record open; that metadata was checked against the supplied snapshot on 2026-10-11. The label is preserved for provenance and is not an independent verification of present open status or mathematical correctness. The catalogue literature assessment is dated 2026-08-22.
\subsection{Short English statement}
Determine whether existing learning theories can explain the reported foundation-model phenomena.
\subsection{Sources}
\begin{itemize}
\item[S1] Ulam AI and the UnsolvedMath Contributors, UnsolvedMath: A Curated Collection of Open Mathematics Problems, record 30006672 (OWR-14299916-010). Catalogue curation and metadata: CC BY 4.0. The source extract is retained; definitions, English summaries and qualifications are editorial additions.
\url{https://huggingface.co/datasets/ulamai/UnsolvedMath}.
\item[S2] Modern and Emerging Phenomena in Machine Learning (2026), Oberwolfach Reports, DOI 10.4171/owr/2026/11, pp. 683--736
\url{https://doi.org/10.4171/owr/2026/11}.
\item Source-specific Oberwolfach Reports licences, including CC BY-SA 4.0 where applicable, are retained; see THIRD\_PARTY\_NOTICES.md.
\end{itemize}
\paragraph{Preparation and review.}
The original source and UnsolvedMath contributors retain attribution for the source extract. Codex (GPT-6; exact serving revision not exposed) assisted with definitions, English summaries, scope notes and formatting. No new human mathematical authorship or independently verified solution is asserted. Permanent source ID 30006672 and code OWR-14299916-010 checked against the complete supplied registry and current repository. Distinct upstream records, including repeated questions and broad or damaged extracts, are retained for complete coverage of the source CS category. A repeated record is not a new independent conjecture. Complete mathematical-equivalence and current-literature review has not been performed. Preparation and final review follow contribution prompts C and D, with source gaps disclosed. Structural and rendering checks do not establish a new theorem or open status.
\label{end:30006672}

\section{Research attempt}
\subsection{Literature, source context and scope}
The full catalogue body, including qualifications and historical source status, is preserved above. The original workshop report was inspected in the locally cached publisher PDF text. This attempt distinguishes the source question from any explicitly restricted model introduced below. Repeated records retain their own source IDs and do not become independent conjectures. Elementary results and attributed literature results are not claimed as new original-target progress.

\subsection{Testing an inference from abrupt measured capability}
The source asks whether fundamentally new principles are needed without specifying the phenomena or a criterion for theoretical novelty. We examine one limited inference: an abrupt benchmark improvement must reflect an abrupt change in the underlying training dynamics. The model below is deliberately elementary and is not an account of a foundation model.
Let a scalar parameter start at $\theta(0)=0$ and follow gradient flow for $L(\theta)=\tfrac12(\theta-1)^2$. Define a thresholded benchmark score $A(t)$ to be one when $\theta(t)\ge1/2$ and zero otherwise.
\begin{proposition}
The parameter trajectory and loss are analytic functions of time, while the benchmark score jumps at $t=\log2$.
\end{proposition}
\begin{proof}
The differential equation is $\dot\theta=1-\theta$, whose solution is $\theta(t)=1-e^{-t}$, as verified by differentiation and the initial value. Substitution gives $L(t)=\tfrac12e^{-2t}$. Both are analytic for all real $t$. The threshold inequality is equivalent to $e^{-t}\le1/2$, or $t\ge\log2$. Thus $A$ is zero before that time and one from that time onward.
\end{proof}
The attempted inference fails because a discontinuous evaluation map can create an apparent transition from a smooth ordinary learning process. This does not establish that real capability changes are measurement artifacts, and it does not exclude distinct mechanisms in realistic architectures. It identifies an ambiguity a proposed explanation must resolve before treating a benchmark jump as evidence for new dynamical principles.
\subsection{Verification and remaining gaps}
The scores immediately before and after $\log2$ were checked, together with the smooth loss values. There is no formal way to assign completion to the source's choice between ``new'' and ``extended'' principles without defining the target phenomena, explanatory standard and allowed mathematical framework. The auxiliary example supplies no empirical foundation-model explanation.

\section{References and source attribution}
The primary Oberwolfach workshop report, \url{https://doi.org/10.4171/owr/2026/11}, was inspected on 11 October 2026. Relevant source passages: 2026/11 pp. 729--730, new principles versus extensions of existing theories. Catalogue and contributor attribution remain in the original Sources section. Mathematical arguments labeled as this attempt are proved above; cited external results remain attributed to their authors and are not independently reproved or claimed as our discoveries.


\paragraph{Integration review.} Codex (GPT-6; exact serving revision not exposed) reviewed the written arguments and expanded layout and source-reference wrappers. Repository metadata and the production web renderer were checked using Python 3.13 and JavaScript with MathJax 3.2.2. This model-assisted review is not independent human proof verification; no PDF compilation or formal proof checking is asserted.

\section{Completion Estimate}
\noindent\textbf{COMPLETION: 0\%}
No substantiated new progress toward the entire original target is claimed. The worked results and failed route are delimited above; a known result, illustrative model, or auxiliary obstruction is not counted as original-target completion.

\section{Final status and attribution}
\textbf{Attempt status: unresolved.} This describes this research attempt, not a new ruling on the source catalogue's status or on external literature claims.
The mathematics and this self-audit were generated by OpenAI \texttt{gpt-6-astra}, reasoning \texttt{ultra}, on 11 October 2026. Codex (GPT-6; exact serving revision not exposed) orchestrated the batch and contributed formatting/repository review. No human mathematical authorship, independent human proof checking, or formal verification is claimed. The historical preparation paragraphs in the reproduced catalogue body concern that body only.

Contribution prompt A was applied to the complete definition and supplied source record, followed by the adversarial checks of prompt B. Local reproduction files are \texttt{batch011/checks.py}, \texttt{computations.json}, and the per-ID audit JSON. They were not public links at preparation time. Numerical checks are finite and do not prove unrestricted results. No LaTeX engine was found locally; compilation was not performed by this research agent.

\end{document}
